% Section 11: phase-space coordinates, Weyl operators, Gaussian unitaries in Weyl form, the Fock-space form of a
% Gaussian unitary (Proposition 5 of Tjoa-Cirac), (6), (8), (11), (30), (31), (A3). GA-BMPS milestone B5.3a,
% Sections 7-8 of its proposal (chapter written 30 Oct 2026). New text for vetting (veteenv).
\providecommand{\cU}{\mathcal U}
\providecommand{\cV}{\mathcal V}
\providecommand{\sU}{\mathsf U}
\providecommand{\sS}{\mathsf S}
\providecommand{\dom}{\operatorname{dom}}
\providecommand{\Sp}{\mathrm{Sp}(2N,\R)}
\providecommand{\Un}{\mathrm{U}(N)}
\providecommand{\Gs}{\Gamma_s(\iota)}
\providecommand{\Rh}{\hat R}
\providecommand{\Qd}{\mathcal Q}
\begin{veteenv}
\setlength{\emergencystretch}{3em}
\section{Gaussian unitaries}\label{sec:gaussunitary}

A Gaussian unitary $U$ acts on the quadratures $\Rh$ by an affine symplectic map, $U^\dagger\Rh U=M\Rh+d$~\cite[(A3)]{tjoacirac}.
Since the quadratures are unbounded, we take the exponentiated form of this relation, an identity between bounded operators:
conjugation by $U$ sends each Weyl operator $W(u)=e^{-i\sigma(u,\Rh)}$ to a phase times $W(M^{-1}u)$. The displacement,
squeezing and passive unitaries of Sections~\ref{sec:secondquant} and~\ref{sec:displacement} are Gaussian in this sense, and
irreducibility (Theorem~\ref{thm:dsq-irreducible}) turns the Bloch--Messiah decomposition of $M$ into a factorization of $U$,
Proposition~5 of~\cite{tjoacirac}.

\medskip
\noindent\textbf{Notation for this section.}
\begin{center}
\begin{tabular}{@{}L{0.3\textwidth}L{0.64\textwidth}@{}}
$\iota=\{1,\ldots,N\}$, $\Gs$ & $N\ge1$ modes, except where a statement allows any finite $\iota$; the bosonic Fock space\\
$u=(q,p)$, $\sigma(u,v)$ & a vector of $\R^{2N}$, positions first; $u^T\Omega v$ (Definition~\ref{def:gu-coords})\\
$\alpha(u)$, $\delta(\alpha)$ & $(q+ip)/\sqrt2\in\C^N$ and its inverse (Definition~\ref{def:gu-coords})\\
$\Rh=(\hat q,\hat p)$, $\Qd$ & quadratures on $\Fock$; $\bigcap_j\dom\bar a_j$ (Definition~\ref{def:gu-coords})\\
$M_\zeta$ & the symplectic matrix of $\sS(\zeta)$ (Definition~\ref{def:gu-coords})\\
$W(u)$; $(M,d)$ & Weyl operator; data of a Gaussian unitary (Definition~\ref{def:gu-weyl})\\
$O_\cU$, $\Sigma(r)$ & orthogonal symplectic matrix of $\cU\in\Un$, squeezing matrix (Definitions~\ref{def:bm-osp},~\ref{def:bm-sigma})\\
$\sU_\cU$, $D(\alpha)$, $\sS(\zeta)$, $R(\theta)$ & passive, displacement, squeezing and rotation operators
(Sections~\ref{sec:secondquant},~\ref{sec:displacement})\\
\end{tabular}
\end{center}
The symplectic matrix of a Gaussian unitary is written $M$, since $\sS$ is the squeezing operator; $U$ is a unitary operator on
$\Gs$ and $\cU,\cV$ are unitary matrices.

\subsection{Phase-space coordinates}\label{sec:gu-coords}

\begin{definition}[Phase-space coordinates and quadratures]\label{def:gu-coords}
\begin{enumerate}
\item[(a)] For $u=(q,p)\in\R^{2N}$, $\alpha(u)=(q+ip)/\sqrt2\in\C^N$; for $\alpha\in\C^N$,
$\delta(\alpha)=\sqrt2\,(\Rea\alpha,\operatorname{Im}\alpha)\in\R^{2N}$, the paper's~(A33). The \emph{symplectic form} is
$\sigma(u,v)=u^T\Omega v=\sum_j(q_jp'_j-p_jq'_j)$ for $v=(q',p')$, with $\Omega=\left(\begin{smallmatrix}0&I\\-I&0\end{smallmatrix}\right)$.
\item[(b)] The \emph{quadratures} on $\Fock$ are $\Rh=(\hat q,\hat p)$, $\hat q_j=(a_j+a_j^\dagger)/\sqrt2$ and
$\hat p_j=i(a_j^\dagger-a_j)/\sqrt2$, and $\Qd=\bigcap_j\dom\bar a_j\subseteq\Gs$.
\item[(c)] For $\zeta\in\C^N$, with $c_j$, $s_j$, $\phi_j$ as in Definition~\ref{def:dsq-squeeze}, $M_\zeta$ is the real $2N\times2N$ matrix
with diagonal $N\times N$ blocks
\[
 M_\zeta=\begin{pmatrix}\diag(c_j-s_j\cos\phi_j)&\diag(-s_j\sin\phi_j)\\ \diag(-s_j\sin\phi_j)&\diag(c_j+s_j\cos\phi_j)\end{pmatrix},
\]
where $\cos\phi_j$ and $\sin\phi_j$ are read as $0$ when $\zeta_j=0$ (then $s_j=0$).
\end{enumerate}
\end{definition}

In the mode-pair order $(q_1,p_1,\ldots,q_N,p_N)$, $M_\zeta$ is the direct sum of the $2\times2$ matrices
$c_j1-s_j\left(\begin{smallmatrix}\cos\phi_j&\sin\phi_j\\\sin\phi_j&-\cos\phi_j\end{smallmatrix}\right)$. Formally
$\alpha(u)\cdot a^\dagger-\overline{\alpha(u)}\cdot a=-i\sigma(u,\Rh)$, so $D(\alpha(u))$ is $e^{-i\sigma(u,\Rh)}$, and the quadratures
satisfy $[\Rh_k,\Rh_l]=i\Omega_{kl}$ on $\Fock$.

\begin{lemma}[Phase-space coordinates]\label{lem:gu-coords}
\begin{enumerate}
\item[(a)] $\alpha(\delta(\alpha))=\alpha$ and $\delta(\alpha(u))=u$; $\alpha$ is additive, odd and injective, and
$\lvert\alpha(u)\rvert^2=\lvert u\rvert^2/2$.
\item[(b)] $\sigma$ is bilinear, $\sum_j\bigl(\alpha(u)_j\overline{\alpha(v)_j}-\overline{\alpha(u)_j}\alpha(v)_j\bigr)=-i\sigma(u,v)$, and if
$\sigma(u,w)=0$ for every $u$, then $w=0$. For $M\in\Sp$, $\sigma(Mu,Mv)=\sigma(u,v)$ and $\sigma(M^{-1}u,v)=\sigma(u,Mv)$.
\item[(c)] $\alpha(O_\cU u)=\cU\alpha(u)$ for $\cU\in\Un$.
\item[(d)] $\alpha(M_\zeta u)=L_\zeta\alpha(u)$, with $L_\zeta$ as in Theorem~\ref{thm:dsq-squeeze}(c); $M_\zeta\in\Sp$, $M_{-\zeta}M_\zeta=1$
and $M_\zeta^{-1}=M_{-\zeta}$.
\end{enumerate}
\end{lemma}

\begin{proof}
\begin{enumerate}
\item[(a)] The maps are mutually inverse real-linear bijections, and $\lvert q_j+ip_j\rvert^2/2=(q_j^2+p_j^2)/2$.
\item[(b)] $\alpha(u)_j\overline{\alpha(v)_j}=(q_j+ip_j)(q'_j-ip'_j)/2$ has imaginary part $(p_jq'_j-q_jp'_j)/2$, and the sum in question
is $2i$ times the sum of these imaginary parts. With $u$ the coordinate vectors, $\sigma(u,w)$ runs through $w_{p_j}$ and $-w_{q_j}$.
By Lemma~\ref{lem:bm-symplectic}, $M^T\Omega M=\Omega$, which is the first identity; the second is the first applied to $M^{-1}u$
and $v$.
\item[(c)] With $\cU=X+iY$, Definition~\ref{def:bm-osp} gives $(Xq-Yp)+i(Yq+Xp)=\cU(q+ip)$.
\item[(d)] In mode $j$, with $e^{i\phi}(q-ip)=(q\cos\phi+p\sin\phi)+i(q\sin\phi-p\cos\phi)$,
\[
 c(q+ip)-e^{i\phi}s(q-ip)=\bigl((c-s\cos\phi)q-s\sin\phi\,p\bigr)+i\bigl(-s\sin\phi\,q+(c+s\cos\phi)p\bigr),
\]
which is $\sqrt2\,\alpha(M_\zeta u)_j$. Each $2\times2$ block has determinant $c^2-s^2=1$, so it is symplectic, and so is the direct
sum. Since $L_{-\zeta}L_\zeta=1$ (Theorem~\ref{thm:dsq-squeeze}(c), with $\zeta$ and $-\zeta$ exchanged) and $\alpha$ is injective,
$M_{-\zeta}M_\zeta=1$.\qedhere
\end{enumerate}
\end{proof}

\subsection{Weyl operators and Gaussian unitaries}\label{sec:gu-weyl}

\begin{definition}[Weyl operator; Gaussian unitary]\label{def:gu-weyl}
The \emph{Weyl operator} is $W(u)=D(\alpha(u))$, $u\in\R^{2N}$. A unitary operator $U$ on $\Gs$ is a \emph{Gaussian unitary with
symplectic matrix $M$ and displacement vector $d\in\R^{2N}$}, or with \emph{data} $(M,d)$, if $M\in\Sp$ and
\[
 U^*W(u)U=e^{-i\sigma(u,d)}\,W(M^{-1}u)\qquad\text{for every }u\in\R^{2N}.
\]
\end{definition}

Formally, $U^*e^{-i\sigma(u,\Rh)}U=e^{-i\sigma(u,M\Rh+d)}$ and $\sigma(u,M\Rh)=\sigma(M^{-1}u,\Rh)$, so the definition is the exponentiated
form of $U^\dagger\Rh U=M\Rh+d$. The class contains products: $\sU_{-I}\sS(r)$ with $r>0$ has the symplectic matrix $-\Sigma(r)$
(Theorem~\ref{thm:gu-implement}), of trace $-2\cosh r<-2$, while the symplectic matrix $e^X$ of a single exponential $e^{-iH}$ with
$H$ quadratic, $X$ real and traceless, has $\operatorname{tr}e^X\ge-2$ (for $X^2=-\det X\cdot1$, the trace is $2\cosh\sqrt{-\det X}$ or
$2\cos\sqrt{\det X}$); that $e^{-iH}$ has a symplectic matrix of this form is \vetq{Theorem~\ref{thm:qh-main}}. So the Gaussian unitaries
are not all of the single-exponential form $e^{i\chi}e^{-iH_G}$ of~\cite[(1)]{tjoacirac}; they form the group generated by these
exponentials, as the paper uses when it composes Gaussian unitaries in~(A34) and~(11). This group is the standard class, generated
by phases, Weyl operators and the metaplectic operators~\cite[\vetx{Ch.~4}]{folland}; by Theorem~\ref{thm:gu-prop5} every Gaussian unitary
is a phase times a product of displacement, squeezing and passive unitaries.

\begin{proposition}[Weyl relation and data]\label{prop:gu-weyl}
\begin{enumerate}
\item[(a)] $W(u)W(v)=e^{-i\sigma(u,v)/2}W(u+v)$, $W(u)^*=W(-u)$, and $\langle0\rvert W(u)\ket0=e^{-\lvert u\rvert^2/4}$.
\item[(b)] If $U_1,U_2$ are Gaussian with data $(M_1,d_1)$, $(M_2,d_2)$, then $U_1U_2$ is Gaussian with data $(M_1M_2,d_1+M_1d_2)$;
if $U$ is Gaussian with data $(M,d)$, then $U^*$ is Gaussian with data $(M^{-1},-M^{-1}d)$; and a unitary $U=\lambda1$ is Gaussian
with data $(1,0)$.
\item[(c)] If $\lambda W(u)=\lambda'W(v)$ with $\lvert\lambda\rvert=\lvert\lambda'\rvert=1$, then $u=v$ and $\lambda=\lambda'$. If $U$ is Gaussian with
data $(M,d)$ and with data $(M',d')$, then $M=M'$ and $d=d'$. If $U$ and $V$ are Gaussian with the same data, then $U=\lambda V$ with
$\lvert\lambda\rvert=1$.
\end{enumerate}
\end{proposition}

\begin{proof}
\begin{enumerate}
\item[(a)] By Theorem~\ref{thm:dsq-disp}(c), Lemma~\ref{lem:gu-coords}(a),(b) and additivity of $\alpha$; $W(u)^*=D(-\alpha(u))=W(-u)$; and
$\langle0\rvert W(u)\ket0=e^{-\lvert\alpha(u)\rvert^2/2}=e^{-\lvert u\rvert^2/4}$ by Theorem~\ref{thm:dsq-disp}(d).
\item[(b)] $(U_1U_2)^*W(u)U_1U_2=e^{-i\sigma(u,d_1)}U_2^*W(M_1^{-1}u)U_2=e^{-i\sigma(u,d_1)-i\sigma(M_1^{-1}u,d_2)}W(M_2^{-1}M_1^{-1}u)$, and
$\sigma(M_1^{-1}u,d_2)=\sigma(u,M_1d_2)$. For $U^*$, conjugate the defining relation at $u=Mv$ by $U$: $UW(v)U^*=e^{i\sigma(Mv,d)}W(Mv)$, and
$\sigma(Mv,d)=\sigma(v,M^{-1}d)$. A scalar unitary commutes with $W(u)$.
\item[(c)] $W(-v)W(u)=e^{i\sigma(v,u)/2}W(u-v)$ by (a), so $\lambda W(u)=\lambda'W(v)$ makes $W(u-v)$ a scalar of modulus $1$; then
$e^{-\lvert u-v\rvert^2/4}=\lvert\langle0\rvert W(u-v)\ket0\rvert=1$ gives $u=v$, and $\lambda=\lambda'$. For two data of $U$, this gives
$M^{-1}u=M'^{-1}u$ for every $u$, so $M=M'$, and $e^{-i\sigma(u,d-d')}=1$ for every $u$. If $\sigma(u,d-d')=t\ne0$ for some $u$, then at
$\pi u/t$ the phase is $e^{-i\pi}=-1$; so $\sigma(\cdot,d-d')=0$ and $d=d'$ by Lemma~\ref{lem:gu-coords}(b). If $U$, $V$ have the same
data, then $U^*D(\alpha)U=V^*D(\alpha)V$ for every $\alpha=\alpha(\delta(\alpha))$, and Corollary~\ref{cor:dsq-irreducible}(a)
applies.\qedhere
\end{enumerate}
\end{proof}

\begin{theorem}[Gaussian unitaries on Fock space;~{\cite[(30), (31), (A17)--(A33)]{tjoacirac}}]\label{thm:gu-implement}
\begin{enumerate}
\item[(a)] $D(\alpha)$ is Gaussian with data $(1,\delta(\alpha))$.
\item[(b)] For every finite set $\iota$, $\cU\in\mathrm U(\iota)$ and $\beta,z\in\C^\iota$: $\sU_\cU^*D(\beta)\sU_\cU=D(\cU^*\beta)$, the paper's (30);
$\sU_\cU D(\beta)=D(\cU\beta)\sU_\cU$; $\sU_\cU\varepsilon(z)=\varepsilon(\cU z)$ and $\sU_\cU\ket0=\ket0$. For $\cU\in\Un$, $\sU_\cU$ is Gaussian with
data $(O_\cU,0)$, the paper's (A17)--(A19).
\item[(c)] $\sS(\zeta)$ is Gaussian with data $(M_\zeta,0)$, and $M_r=\Sigma(r)$ for every $r\in\R^N$, the paper's (A27)--(A29).
\item[(d)] For $\cU\in\Un$, a mode $j$ and $\zeta\in\C$, $\sU_\cU^*\sS(\zeta e_j)\sU_\cU$ is Gaussian with data
$(O_\cU^TM_{\zeta e_j}O_\cU,0)$. With $b_j=\sum_k\cU_{jk}a_k$ and $b_j^\dagger=\sum_k\overline{\cU_{jk}}a_k^\dagger$: if $\varphi\in\Gs$ is finitely
supported and $G\in\Gs$ satisfies $\jmath G=\bigl(\tfrac12\bar\zeta b_j^2-\tfrac12\zeta b_j^{\dagger2}\bigr)\jmath\varphi$, then
$t\mapsto\sU_\cU^*\sS(t\zeta e_j)\sU_\cU\varphi$ is differentiable at $t=0$ with derivative $G$: the paper's (31).
\end{enumerate}
\end{theorem}

\begin{proof}
\begin{enumerate}
\item[(a)] Applying the Weyl relation twice, $D(-\alpha)D(\beta)D(\alpha)=e^{\langle\alpha,\beta\rangle-\langle\beta,\alpha\rangle}D(\beta)$. For
$\beta=\alpha(u)$ and $\alpha=\alpha(d)$, $d=\delta(\alpha)$, Lemma~\ref{lem:gu-coords}(b) gives
$\langle\alpha(d),\alpha(u)\rangle-\langle\alpha(u),\alpha(d)\rangle=-i\sigma(u,d)$.
\item[(b)] By Lemma~\ref{lem:dsq-expvec}(b) and Theorem~\ref{thm:dsq-disp}(b), $\sU_{\cU^*}D(\beta)\sU_\cU\varepsilon(z)=
e^{-\lvert\beta\rvert^2/2-\langle\beta,\cU z\rangle}\varepsilon(z+\cU^*\beta)$, which is $D(\cU^*\beta)\varepsilon(z)$ because
$\lvert\cU^*\beta\rvert=\lvert\beta\rvert$ and $\langle\beta,\cU z\rangle=\langle\cU^*\beta,z\rangle$; Theorem~\ref{thm:dsq-total} and
$\sU_\cU^*=\sU_{\cU^*}$ (Theorem~\ref{thm:sq-unitary}(b)) give (30). Replacing $\beta$ by $\cU\beta$ and multiplying by $\sU_\cU$ gives the
second identity, and $\sU_\cU\ket0=\ket0$ is $z=0$. For $\cU\in\Un$, $\cU^*\alpha(u)=\alpha(O_{\cU^*}u)=\alpha(O_\cU^{-1}u)$ by
Lemma~\ref{lem:gu-coords}(c) and Theorem~\ref{thm:bm-unitary}(a), and $O_\cU\in\Sp$ by Theorem~\ref{thm:bm-unitary}(b).
\item[(c)] By Theorem~\ref{thm:dsq-squeeze}(c) and Lemma~\ref{lem:gu-coords}(d), $\sS(\zeta)^*W(u)\sS(\zeta)=D(L_{-\zeta}\alpha(u))=
W(M_{-\zeta}u)=W(M_\zeta^{-1}u)$. For real $r_j>0$, $\phi_j=0$ and the blocks are $c_j\mp s_j=e^{\mp r_j}$; for $r_j<0$, $\phi_j=\pi$,
$s_j=\sinh\lvert r_j\rvert$ and the blocks are $c_j\pm s_j=e^{\mp r_j}$; for $r_j=0$ they are $1$.
\item[(d)] The data follow from (b), (c) and Proposition~\ref{prop:gu-weyl}(b), with $O_\cU^{-1}=O_\cU^T$. The vector $\sU_\cU\varphi$ is
finitely supported (Lemma~\ref{lem:sq-sectors}(c)), so by Theorem~\ref{thm:dsq-squeeze-group}(c) $t\mapsto\sS(t\zeta e_j)\sU_\cU\varphi$
has at $0$ the derivative $G_0$ with $\jmath G_0=(\frac12\bar\zeta a_j^2-\frac12\zeta a_j^{\dagger2})\Gamma(\cU)\jmath\varphi$; $G_0$ exists by
Lemma~\ref{lem:dsq-above}. The bounded operator $\sU_\cU^*$ maps it to the derivative of the conjugated group, and
$\jmath\sU_\cU^*G_0=\Gamma(\cU^*)(\frac12\bar\zeta a_j^2-\frac12\zeta a_j^{\dagger2})\Gamma(\cU)\jmath\varphi=\jmath G$ by
Theorems~\ref{thm:sq-unitary}(c) and~\ref{thm:sq-conj}(a).\qedhere
\end{enumerate}
\end{proof}

By (b), the product form of the paper's (30), $\sU_\cU^*D(\beta_0e_j)\sU_\cU=\prod_kD(\overline{\cU_{jk}}\beta_0e_k)$, is the product formula of
Theorem~\ref{thm:dsq-disp}(d) for $\cU^*(\beta_0e_j)$, and $R(\theta)$ is Gaussian with data $(O_{\diag(e^{-i\theta_j})},0)$. Conjugating
the group law and the strong continuity of Theorem~\ref{thm:dsq-squeeze-group}(a),(b) by the fixed unitary $\sU_\cU$ shows that
$t\mapsto\sU_\cU^*\sS(t\zeta e_j)\sU_\cU$ is a strongly continuous one-parameter unitary group; by (d) its generator acts on finitely
supported vectors as the squeezing generator in the modes $b_j$.

\subsection{The Fock-space form of a Gaussian unitary}\label{sec:gu-prop5}

\begin{theorem}[{\cite[Prop.~5, (A12)]{tjoacirac}}, for Gaussian unitaries in the sense of Definition~\ref{def:gu-weyl}]\label{thm:gu-prop5}
Let $U$ be a Gaussian unitary with data $(M,d)$.
\begin{enumerate}
\item[(a)] If $M=O_\cU\Sigma(r)O_\cV$ with $\cU,\cV\in\Un$ and $r\in\R^N$, then $U=\lambda\,\sU_\cU D(\alpha)\sS(r)\sU_\cV$ for some
$\lambda$ with $\lvert\lambda\rvert=1$, where $\alpha=\alpha(O_\cU^Td)$, that is $\delta(\alpha)=O_\cU^Td$. In particular there are $\lambda$ with
$\lvert\lambda\rvert=1$, $\cU,\cV\in\Un$, $\alpha\in\C^N$ and $r\in\R^N$ with $r\ge0$ such that $U=\lambda\,\sU_\cU D(\alpha)\sS(r)\sU_\cV$.
\item[(b)] For $\cU,\cV\in\Un$ and $\alpha,\zeta\in\C^N$, the unitary $\sU_\cU D(\alpha)\sS(\zeta)\sU_\cV$ is Gaussian with the
data $(O_\cU M_\zeta O_\cV,\,O_\cU\delta(\alpha))$, the paper's (A35).
\item[(c)] Let $\iota$ be a finite set, $\zeta\in\C^\iota$ and $\theta\in\R^\iota$ with $\lvert\zeta_j\rvert e^{-2i\theta_j}=\zeta_j$ for
every $j$, and write $\Theta_\theta=\diag(e^{-i\theta_j})$, so that $R(\theta)=\sU_{\Theta_\theta}$. Then $\sS(\zeta)=\sU_{\Theta_\theta}\,\sS(\lvert\zeta\rvert)\,
\sU_{\Theta_\theta}^*$, where $\lvert\zeta\rvert=(\lvert\zeta_j\rvert)_j$. For $\iota=\{1,\ldots,N\}$,
$\sU_\cU D(\alpha)\sS(\zeta)\sU_\cV=\sU_{\cU\Theta_\theta}D(\Theta_\theta^*\alpha)\sS(\lvert\zeta\rvert)\sU_{\Theta_\theta^*\cV}$.
\end{enumerate}
\end{theorem}

The admissible $\theta$ in (c) are $\theta_j=-\phi_j/2+k_j\pi$, $k_j\in\mathbb Z$, where $\zeta_j\ne0$, and arbitrary where
$\zeta_j=0$; the choice $\theta_j=-\phi_j/2$ gives
$\Theta_\theta=\diag(e^{i\phi_j/2})$, the paper's absorption of the squeezing phases into the passive factors~(A38). The numbers
$r_j\ge0$ in (a) are determined by $M$ up to order (Proposition~\ref{prop:bm-singular}(b)), and $\lambda$ is determined by $U$ and the
factors, since the product of the factors is unitary.

\begin{proof}
\begin{enumerate}
\item[(b)] By Theorem~\ref{thm:gu-implement} the factors have data $(O_\cU,0)$, $(1,\delta(\alpha))$, $(M_\zeta,0)$, $(O_\cV,0)$, and by
Proposition~\ref{prop:gu-weyl}(b) the product has data $(O_\cU M_\zeta O_\cV,\,O_\cU\delta(\alpha))$.
\item[(a)] By (b) with $\zeta=r$, $M_r=\Sigma(r)$, the product $\sU_\cU D(\alpha)\sS(r)\sU_\cV$ has data $(M,O_\cU\delta(\alpha))=(M,d)$,
the paper's (A36); Proposition~\ref{prop:gu-weyl}(c) gives $U=\lambda\,\sU_\cU D(\alpha)\sS(r)\sU_\cV$, the paper's (A37). For the existence
form, Theorem~\ref{thm:bm-decomp} writes $M=O\Sigma(r)O'$ with $O,O'$ orthogonal symplectic and $r\ge0$, and
Theorem~\ref{thm:bm-unitary}(b) writes $O=O_\cU$, $O'=O_\cV$.
\item[(c)] By Theorem~\ref{thm:gu-implement}(b) and Theorem~\ref{thm:dsq-squeeze}(c), the unitary $\sU_{\Theta}\sS(\lvert\zeta\rvert)\sU_{\Theta}^*$,
$\Theta=\Theta_\theta$, conjugates $D(\alpha)$ to $D(\Theta L_{-\lvert\zeta\rvert}\Theta^*\alpha)$, and in mode $j$
$e^{-i\theta_j}(c_je^{i\theta_j}\alpha_j+s_j\overline{e^{i\theta_j}\alpha_j})=c_j\alpha_j+e^{-2i\theta_j}s_j\bar\alpha_j=(L_{-\zeta}\alpha)_j$, since
$e^{-2i\theta_j}=e^{i\phi_j}$ when $\zeta_j\ne0$ and $s_j=0$ otherwise. So it has the conjugation action of $\sS(\zeta)$. Its vacuum
expectation is $\langle\sU_\Theta^*\ket0,\sS(\lvert\zeta\rvert)\sU_\Theta^*\ket0\rangle=\langle0\rvert\sS(\lvert\zeta\rvert)\ket0=(\prod_jc_j)^{-1/2}$,
which is $\langle0\rvert\sS(\zeta)\ket0\ne0$ (Theorem~\ref{thm:dsq-squeeze}(a)). Corollary~\ref{cor:dsq-irreducible}(a) gives the first identity.
The second follows from it, from $D(\alpha)\sU_\Theta=\sU_\Theta D(\Theta^*\alpha)$ (Theorem~\ref{thm:gu-implement}(b)) and from
$\sU_\cU\sU_\Theta=\sU_{\cU\Theta}$.\qedhere
\end{enumerate}
\end{proof}

\begin{corollary}[The vacuum, one mode, and the quadratures]\label{cor:gu-vacuum}
\begin{enumerate}
\item[(a)] For one mode and $\alpha,\zeta\in\C$, $D(\alpha)\sS(\zeta)\ket0=c^{-1/2}e^{-\lvert\alpha\rvert^2/2-\tau\bar\alpha^2/2}\ket{\kappa,\ell}$ with
$\kappa=-\tau/2$ and $\ell=\alpha+\tau\bar\alpha=\alpha-2\kappa\bar\alpha$, the paper's (8) with its constant. For $N$ modes and
$\alpha,\zeta\in\C^N$,
\[
 D(\alpha)\sS(\zeta)\ket0=\Bigl(\prod_jc_j\Bigr)^{-1/2}e^{-\lvert\alpha\rvert^2/2-\sum_j\tau_j\bar\alpha_j^2/2}\bigotimes_j\ket{-\tau_j/2,\,\alpha_j+\tau_j\bar\alpha_j}.
\]
If $U=\lambda\,\sU_\cU D(\alpha)\sS(r)\sU_\cV$ with $\lambda\in\C$, $\cU,\cV\in\Un$, $\alpha\in\C^N$ and $r\in\R^N$, then
$U\ket0=\lambda\,\sU_\cU\bigotimes_j\bigl(D(\alpha_j)\sS(r_j)\ket0\bigr)$, the paper's (11) and (A40)--(A42). Every Gaussian unitary $U$
satisfies this for some $\lambda$ with $\lvert\lambda\rvert=1$, $\cU$, $\cV$, $\alpha$ and $r\ge0$.
\item[(b)] For one mode, every Gaussian unitary is $\lambda D(\alpha)\sS(\zeta)R(\theta)$ for some $\lambda,\alpha,\zeta\in\C$ and $\theta\in\R$
with $\lvert\lambda\rvert=1$, the paper's (6). Moreover $R(\theta)R(\theta')=R(\theta+\theta')$ on any number of modes, and every unitary
$1\times1$ matrix is $e^{-i\theta}$ for some real $\theta$.
\item[(c)] Let $U$ be Gaussian with data $(M,d)$. Then $U\psi\in\Qd$ if and only if $\psi\in\Qd$, so $U$ maps $\Qd$ onto itself. If
$\psi\in\Qd$ and $\xi\in\Gs$ satisfies $\jmath\xi=\sum_lM_{kl}\Rh_l\jmath\psi+d_k\jmath\psi$, then $\Rh_k\jmath U\psi=\jmath U\xi$; that is,
$U^*\Rh_kU=\sum_lM_{kl}\Rh_l+d_k$ on $\Qd$, the paper's (A3). Such a $\xi$ exists for every $\psi\in\Qd$, all real $M$, $d$ and every $k$.
\end{enumerate}
\end{corollary}

\begin{proof}
\begin{enumerate}
\item[(a)] We compare inner products with $\varepsilon(w)$, $w\in\C^N$. By Theorem~\ref{thm:dsq-disp}(b),
$D(-\alpha)\varepsilon(w)=e^{-\lvert\alpha\rvert^2/2+\langle\alpha,w\rangle}\varepsilon(w-\alpha)$, so by Theorem~\ref{thm:dsq-squeeze}(a) at $z=0$,
\[
 \langle\varepsilon(w),D(\alpha)\sS(\zeta)\ket0\rangle=\Bigl(\prod_jc_j\Bigr)^{-1/2}
 \exp\sum_j\Bigl(-\tfrac12\lvert\alpha_j\rvert^2+\alpha_j\bar w_j-\tfrac12\tau_j(\bar w_j-\bar\alpha_j)^2\Bigr).
\]
By the remark after Lemma~\ref{lem:dsq-mehler}, the right side of the claim gives $(\prod_jc_j)^{-1/2}$ times
$\exp\sum_j(-\tfrac12\lvert\alpha_j\rvert^2-\tfrac12\tau_j\bar\alpha_j^2-\tfrac12\tau_j\bar w_j^2+(\alpha_j+\tau_j\bar\alpha_j)\bar w_j)$, which is
equal. Both vectors lie in
$\jmath(\Gs)$, the product vector because $\lvert\tau_j/2\rvert<\frac12$ (proof of Lemma~\ref{lem:dsq-lower}), so they are equal by
Theorem~\ref{thm:dsq-total}; one mode is the case $N=1$. For the given $U$, $\sU_\cV\ket0=\ket0$, and the several-mode formula is the
product vector of the one-mode formulas, since $\bigotimes_j(x_j\chi_j)=(\prod_jx_j)\bigotimes_j\chi_j$ for numbers $x_j$. The existence form is
Theorem~\ref{thm:gu-prop5}(a).
\item[(b)] A unitary $1\times1$ matrix is a number of modulus $1$, so $e^{-i\theta}$, and $R(\theta)R(\theta')=R(\theta+\theta')$ because
$\Theta_\theta\Theta_{\theta'}=\Theta_{\theta+\theta'}$ and $\cU\mapsto\sU_\cU$ is a homomorphism. By Theorem~\ref{thm:gu-prop5}(a),
$U=\lambda R(\theta_1)D(\alpha)\sS(r)R(\theta_2)$ with $r\ge0$. By Theorem~\ref{thm:gu-implement}(b), $R(\theta_1)D(\alpha)=D(e^{-i\theta_1}\alpha)R(\theta_1)$.
By Theorem~\ref{thm:gu-prop5}(c) with $\zeta=re^{-2i\theta_1}$, which satisfies $\lvert\zeta\rvert e^{-2i\theta_1}=\zeta$ for every real $\theta_1$
because $r\ge0$, $R(\theta_1)\sS(r)=\sS(re^{-2i\theta_1})R(\theta_1)$. Hence $U=\lambda D(e^{-i\theta_1}\alpha)\sS(re^{-2i\theta_1})R(\theta_1+\theta_2)$.
\item[(c)] Say that a unitary $V$ \emph{implements} $(M,d)$ on $\Qd$ if $V$ maps $\Qd$ onto $\Qd$ and $\Rh_k\jmath V\psi=\jmath V\xi$ whenever
$\psi\in\Qd$ and $\jmath\xi=\sum_lM_{kl}\Rh_l\jmath\psi+d_k\jmath\psi$. A phase implements $(1,0)$. By Theorem~\ref{thm:dsq-squeeze-intertwine},
$\Qd=\bigcap_j\dom\overline{a_j^\dagger}$ as well, so for $\psi\in\Qd$ every $\Rh_l\jmath\psi$ lies in $\jmath(\Gs)$, and $\xi$ exists. By
Theorem~\ref{thm:dsq-disp-intertwine}, with $(\alpha_j+\bar\alpha_j)/\sqrt2=\sqrt2\Rea\alpha_j$ and $i(\bar\alpha_j-\alpha_j)/\sqrt2=\sqrt2\operatorname{Im}\alpha_j$,
$D(\alpha)$ implements $(1,\delta(\alpha))$. By Theorem~\ref{thm:dsq-squeeze-intertwine}, $\sS(\zeta)$ implements the matrix of
$a\mapsto c\,a-e^{i\phi}s\,a^\dagger$ on $(\hat q,\hat p)$, which is $M_\zeta$ by the computation in the proof of Lemma~\ref{lem:gu-coords}(d).
The passive unitary maps $\Qd$ into $\Qd$, since $a_j\Gamma(\cU)\jmath\psi=\Gamma(\cU)\sum_k\cU_{jk}a_k\jmath\psi$
(Theorem~\ref{thm:sq-intertwine}) and $\Gamma(\cU)$ preserves $\jmath(\Gs)$; with $\sU_{\cU^*}$ it maps onto. By
Theorem~\ref{thm:sq-unitary}(c), $\Gamma(\cU^*)\Rh\Gamma(\cU)=O_\cU\Rh$ on $\Fock$, since $(X+iY)(\hat q+i\hat p)=(X\hat q-Y\hat p)+i(Y\hat q+X\hat p)$;
so $\sU_\cU$ implements $(O_\cU,0)$. If $V_1$ implements $(M_1,d_1)$ and $V_2$ implements $(M_2,d_2)$, then $V_1V_2$ implements
$(M_1M_2,d_1+M_1d_2)$: with $\jmath\xi_l=\sum_m(M_2)_{lm}\Rh_m\jmath\psi+(d_2)_l\jmath\psi$,
\[
 \Rh_k\jmath V_1V_2\psi=\jmath V_1\Bigl(\sum_l(M_1)_{kl}V_2\xi_l+(d_1)_kV_2\psi\Bigr)=\jmath V_1V_2\Bigl(\sum_l(M_1)_{kl}\xi_l+(d_1)_k\psi\Bigr),
\]
and $\jmath$ of the last bracket is $\sum_m(M_1M_2)_{km}\Rh_m\jmath\psi+(d_1+M_1d_2)_k\jmath\psi$. By Theorem~\ref{thm:gu-prop5}(a), $U$ is a phase
times $\sU_\cU D(\alpha)\sS(r)\sU_\cV$, which implements $(O_\cU\Sigma(r)O_\cV,O_\cU\delta(\alpha))=(M,d)$, the order of~(A35).\qedhere
\end{enumerate}
\end{proof}

In (b) the sign of $e^{-i\theta_1}$ needs no choice: Theorem~\ref{thm:gu-prop5}(c) holds for every admissible $\theta$, not only for
$\theta_j=-\phi_j/2$. Equivalently, $\sU_{-I}$ commutes with $\sS(r)$, since both preserve the parity of the particle number.

\subsection*{Notes and references}

Tjoa and Cirac~\cite{tjoacirac}, Appendix~A, lift the factors of the Bloch--Messiah decomposition $S=O_U\Sigma O_W$~(A11) to Fock space
in (A17)--(A33), compute the action of $\tilde U_G=\mathsf UD(\alpha)S(r)\mathsf W$ on the quadratures~(A35), choose $\delta=O_U^Td$~(A36),
and conclude $U_G=e^{i\chi}\tilde U_G$~(A37) because $U_G\tilde U_G^\dagger$ commutes with every canonical operator and the Fock
representation is irreducible; they absorb the squeezing phases into the passive factors~(A38) and obtain~(11) in (A40)--(A42).
We prove their Proposition~5 for Gaussian unitaries in the Weyl form of Definition~\ref{def:gu-weyl}, with irreducibility proved in
Theorem~\ref{thm:dsq-irreducible}. These unitaries satisfy (A3) on $\Qd$ (Corollary~\ref{cor:gu-vacuum}(c)); the converse, from
(A3) with the paper's unspecified domains to the Weyl form, needs input on self-adjointness and domains and is not proved here.
The overlaps of one-mode Gaussian vectors in Corollary~\ref{cor:gu-vacuum}(a) come from Lemma~\ref{lem:dsq-mehler}, the
one-mode case of Mehler's formula, which the companion notes on Gaussian-augmented matrix product states use for matrix
arguments. The metaplectic representation and the group of Gaussian unitaries are treated in~\cite{folland}.
\end{veteenv}
